\section{A hierarchy of geometric statistical statements}
\label{sec:geometric_eth_ansatz}

The bundle construction specifies observables but not a probability law.  A
statistical statement requires an ensemble and a limiting procedure.  We
therefore consider a sequence indexed by $L$, where $L$ may denote particle
number, volume, fiber rank, or another explicitly stated control.  Within
each $L$, all exact symmetries and invariant blocks are resolved before
samples are pooled.  A sample may be a disorder realization, a base point in
a stationary parameter ensemble, or a Hamiltonian drawn from a specified
population.  These choices are not interchangeable and must be recorded as
part of the definition.

Let $a$ label an accessible tangent direction and write
$x_a=\operatorname{vec}(X_a)$ in a local frame.  We subtract the population
mean and apply a deterministic normalization chosen without using the
higher-order outcomes under test:
\begin{equation}
 z_a=W_L\left(x_a-\mathbb E x_a\right).
 \label{eq:normalized_response}
\end{equation}
Here $W_L$ may remove a scalar variance or whiten a fixed symmetry-resolved
support.  A frame change acts linearly on $z_a$; the tensors below transform
covariantly, while the final statements are made for their invariant
contractions.  Outcome-dependent whitening, post-selection of tangents, and
refitting $W_L$ after inspecting fourth-order data are excluded.

Complex response variables require two second-moment tensors.  Suppressing
the tangent label, define
\begin{equation}
 C_{ij}=\mathbb E[z_i z_j^*],
 \qquad
 \Pi_{ij}=\mathbb E[z_i z_j].
 \label{eq:covariance_pseudocovariance}
\end{equation}
We use $\Pi$ for pseudocovariance in the manuscript to avoid collision with
the spectral projector $P$.  A proper complex population has $\Pi=0$, but
properness must be checked rather than imposed.  For a centered complex
Gaussian vector the representative fourth moment is
\begin{equation}
 \mathbb E[z_i z_j^*z_k z_l^*]
 =C_{ij}C_{kl}+C_{il}C_{kj}+\Pi_{ik}\Pi_{jl}^* .
 \label{eq:three_wick_pairings}
\end{equation}
The right-hand side contains all three Wick pairings.  Subtracting only the
two covariance pairings defines an incomplete null whenever $\Pi\ne0$.
Equivalently, one may realify $z$ into $(\operatorname{Re}z,
\operatorname{Im}z)$ and use its full real covariance.

We next define three levels of a possible Geometric-ETH statement.  They are
definitions of claims, not results of the present calculation.

\paragraph{Weak geometric closure.}
A sequence has weak geometric closure for a stated set of tangents and
invariant observables if: (i) the centered, normalized two-point tensors
$(C_L,\Pi_L)$ have deterministic limits on their resolved support; and (ii)
the selected quadratic gauge invariants, such as
$D^{-1}\operatorname{Tr}\widetilde F_{\mu\nu}^2$ or contractions of
$\mathcal Q_{\mu\nu}$, concentrate around the values determined by those
limits.  This definition makes no Gaussian assertion and places no condition
on connected cumulants of order four and above.  It is weaker than ordinary
ETH because it contains neither the smooth diagonal function
$\overline O(E)$ nor a thermalization statement.

\paragraph{Strong geometric closure.}
A sequence has strong geometric closure if it has weak closure and, after
complete whitening on every symmetry-resolved support, every fixed
finite-dimensional real-linear projection has a Gaussian limit:
\begin{equation}
 \kappa_k\!\left(
 f_1(z_L),\ldots,f_k(z_L)
 \right)\longrightarrow0,
 \qquad k\ge3,
 \label{eq:strong_geometric_closure}
\end{equation}
for all fixed test functionals for which the required moments exist.  The
limiting complex Gaussian is specified by both $C$ and $\Pi$.  We additionally
require the accessible response algebra to have no unresolved invariant
block in the fiber.  This irreducibility condition is necessary for a
full-matrix law, but it does not imply
Eq.~\eqref{eq:strong_geometric_closure}.  Strong geometric closure remains a
statement about parameter-space response, not about real-time thermalization.

\paragraph{Deformed geometric closure.}
A sequence has deformed geometric closure if the complete second moments
stabilize but specified connected tensors retain a reproducible structure
after whitening.  Examples include a fixed improper-complex component,
channel-dependent fourth cumulants, locality blocks, unequal gap weights in
Eq.~\eqref{eq:kato_response}, or a nonzero limiting connected tensor.  The
deformation must be stated before comparison and carried by measured tensors
or a specified probability law.  Calling any non-Gaussian residual
``deformed'' without such a prediction would be tautological.  A reducible
symmetry mixture is also not a deformation; the blocks must first be
separated.

These definitions distinguish covariance closure, Gaussian closure, and
structured non-Gaussian closure.  The distinction matters because a
Wishart-like QGT can arise whenever $X_a$ is approximately Gaussian, whereas
curvature is a difference of correlated bilinears and can preserve memory of
$\Pi$, unequal channel weights, or local constraints.  Agreement of an
eigenvalue histogram with a familiar random-matrix curve tests only the
particular marginal shown.  Strong closure requires the complete covariance
and higher connected tensors, not a single spectral statistic.

The model used for the exact result below is more restrictive than any of
these definitions.  It assumes that a normalized response can be decomposed
as
\begin{equation}
 Z=\sum_{\alpha=1}^{n}w_\alpha Y_\alpha,
 \qquad \mathbb E Y_\alpha=0,
 \label{eq:random_channel_model}
\end{equation}
where the $Y_\alpha$ are independent complex random matrices in one common
response space and the real weights $w_\alpha$ are deterministic.  We call
Eq.~\eqref{eq:random_channel_model} the \emph{independent-channel model}.
It is not a consequence of exact degeneracy, a spectral gap, or an
irreducible response algebra.  Establishing it for a physical family requires
identifying the independent units and showing that mixed connected
cumulants are absent or controlled.

The distinction between definition and model fixes the logic of the paper.
Section~\ref{sec:channel_theorem} derives an exact property of
Eq.~\eqref{eq:random_channel_model}.  The numerical oracle checks the
implementation on synthetic channels drawn from the assumed population.
The chiral calculation then asks prospectively whether a new protected model
obeys the resulting fixed prediction.  It does not: the random tangent class
misses the sealed validation band.  Consequently neither weak, strong, nor
deformed geometric closure is established here for a general physical model
class.  The failed test instead identifies independence and channel
stationarity as substantive hypotheses rather than consequences of a large
degenerate subspace.

For later reference, we reserve \emph{Geometric ETH} for the conjecture that
an explicitly specified family has at least weak closure and, after all
protection-induced structures are accounted for, approaches either the
strong or a preregistered deformed class.  This conjecture can be true for one
protection and tangent ensemble and false for another.  No asymptotic,
universal, gravitational, or thermalization form of the conjecture follows
from the finite-channel theorem.
